% ==========================================
% Built into the book on 2026-10-24 by Claude (claude-opus-5-5), from the pipeline's step-4
% output below. The pipeline named a lecturer ("Alexander Gorodnik") who did not give this
% lecture; the name is removed (the user confirmed it is wrong, and the lecturer refers back to
% "Theorem B that I said" yesterday). Further: week-tagged labels and \recap lines, recall
% anchors and cross-references, the math fixes of a review, and the end-of-lecture checkpoint
% resolved (the proof is continued in the next lecture; its completion is the aiNote).
% ==========================================
% ==========================================
% Step 4/4 - Final refinement (refinement step 3)
% Output: step4-10-21-2026-wednesday-week6-Algebra_I-offset-final.tex
% Input: step3-10-21-2026-wednesday-week6-Algebra_I-offset-speech_refined.tex
% Model: gemini-3.8-flash
% Temperature: 0.35
% TopP: 0.9
% TopK: 10
% MaxOutputTokens: 65535
% ThinkingBudget: 4096
% ThinkingLevel: MEDIUM
% Processed on: 2026-10-24 14:19:36
% ------------------------------------------
% Token Usage (1 Request(s)):
%   - Total Prompt Tokens : 56'561 (gesamter Prompt, inkl. Cache)
%   - Cached Context      : 0 (aus Google Context-Cache, rabattiert)
%   - Fresh Input Tokens  : 56'561 (neu verrechneter Input)
%   - Output Tokens       : 23'646 (generiertes LaTeX, ohne Thinking)
%   - Thinking Tokens     : 23'673 (Denk-Tokens, als Output verrechnet)
% ==========================================
% ==========================================
% Step 3/4 - Speech refinement (refinement step 2)
% Output: step3-10-21-2026-wednesday-week6-Algebra_I-offset-speech_refined.tex
% Input: step2-10-21-2026-wednesday-week6-Algebra_I-offset-merged.tex
% Model: gemini-3.8-flash
% Temperature: 0.35
% TopP: 0.9
% TopK: 10
% MaxOutputTokens: 65535
% ThinkingBudget: 4096
% ThinkingLevel: MEDIUM
% Processed on: 2026-10-24 14:09:40
% ------------------------------------------
% Token Usage (1 Request(s)):
%   - Total Prompt Tokens : 151'195 (gesamter Prompt, inkl. Cache)
%   - Cached Context      : 0 (aus Google Context-Cache, rabattiert)
%   - Fresh Input Tokens  : 151'195 (neu verrechneter Input)
%   - Output Tokens       : 22'279 (generiertes LaTeX, ohne Thinking)
%   - Thinking Tokens     : 11'456 (Denk-Tokens, als Output verrechnet)
% ==========================================
% ==========================================
% Step 2/4 - Merge & timestamp alignment (refinement step 1)
% Output: step2-10-21-2026-wednesday-week6-Algebra_I-offset-merged.tex
% Input: step1-10-21-2026-wednesday-week6-Algebra_I-all-offset.tex
% Model: gemini-3.8-flash
% Temperature: 0.4
% TopP: 0.95
% TopK: 40
% MaxOutputTokens: 65535
% ThinkingBudget: 32765
% ThinkingLevel: MEDIUM
% Processed on: 2026-10-24 14:02:58
% ------------------------------------------
% Token Usage (1 Request(s)):
%   - Total Prompt Tokens : 108'826 (gesamter Prompt, inkl. Cache)
%   - Cached Context      : 0 (aus Google Context-Cache, rabattiert)
%   - Fresh Input Tokens  : 108'826 (neu verrechneter Input)
%   - Output Tokens       : 23'182 (generiertes LaTeX, ohne Thinking)
%   - Thinking Tokens     : 11'693 (Denk-Tokens, als Output verrechnet)
% ==========================================
% ==========================================
% Step 1/4 - Transcription, 2 part(s) combined
% Output: step1-10-21-2026-wednesday-week6-Algebra_I-all-offset.tex
% Source video: 10-21-2026-wednesday-week6-Algebra_I.mp4
% Model: gemini-3.8-flash
% Temperature: 0.36
% TopP: 0.95
% TopK: 40
% MaxOutputTokens: 65535
% ThinkingBudget: 32765
% ThinkingLevel: HIGH
% Processed on: 2026-10-24 13:53:55
% ------------------------------------------
% Parts:
%   - Part 1: 10-21-2026-wednesday-week6-Algebra_I-part1.mp4 | start 00:00:00 | model gemini-3.8-flash | processed 2026-10-24 13:10:47 | prompt 357'493 | output 11'656 | thinking 24'529
%   - Part 2: 10-21-2026-wednesday-week6-Algebra_I-part2.mp4 | start 00:19:22 | model gemini-3.8-flash | processed 2026-10-24 13:53:55 | prompt 741'314 | output 13'666 | thinking 78'020
% ------------------------------------------
% Token Usage of this run (3 Request(s)):
%   - Total Prompt Tokens : 1'098'807 (gesamter Prompt, inkl. Cache)
%   - Cached Context      : 262'427 (aus Google Context-Cache, rabattiert)
%   - Fresh Input Tokens  : 836'380 (neu verrechneter Input)
%   - Output Tokens       : 25'322 (generiertes LaTeX, ohne Thinking)
%   - Thinking Tokens     : 102'549 (Denk-Tokens, als Output verrechnet)
% ==========================================

% Primary Language: English
% End of the video: 00:42:45

\lecturechapter[6]{Wednesday}{21 Oct}{21 October 2026}{Gauss's Theorem on Polynomial Rings over UFDs}

\begin{overview}[Lecture Context \& Content Overview]
This lecture proves Gauss's Theorem on polynomial rings (Theorem~B of the previous lecture, \cref{thm:ufd-polynomial-ring}): if $R$ is a unique factorization domain (UFD), then the polynomial ring $R[x]$ is also a unique factorization domain. The last case of the proof is cut off by the end of the hour and finished in the next lecture; its completion is recorded at the end of this chapter.

The lecture is organized into two primary components:
\begin{itemize}
    \item \textbf{Foundational Preliminaries:}
    \begin{itemize}
        \item Review of the equivalence between irreducible elements and prime elements in any UFD.
        \item Definition and computation of the \newterm{content} $c(f)$ of a non-zero polynomial, uniquely determined up to associates in $R$.
        \item Canonical decomposition $f = c(f) \cdot f^*$ into content and a \newterm{primitive} polynomial $f^*$.
        \item Proof of \emph{Gauss's Lemma}: the product of primitive polynomials is primitive ($c(f) \sim 1 \land c(g) \sim 1 \implies c(fg) \sim 1$).
        \item Divisibility property of primitive polynomials: for primitive $g \in R[x]$, $g \mid bf \implies g \mid f$ for any non-zero scalar $b \in R$.
    \end{itemize}
    \item \textbf{Proof of Gauss's Theorem (Theorem B):}
    \begin{itemize}
        \item \textbf{Step 1 (Existence):} Existence of irreducible factorizations in $R[x]$ via induction on $\deg(f)$, using primitivity to rule out non-unit degree-zero factors.
        \item \textbf{Step 2 (Uniqueness):} Uniqueness via proving that every irreducible element in $R[x]$ is prime:
        \begin{itemize}
            \item \emph{Case 1 ($\deg(p) = 0$):} Constant irreducibles in $R$ remain prime in $R[x]$ via coefficient divisibility and Gauss's Lemma.
            \item \emph{Case 2 ($\deg(p) > 0$):} Positive-degree irreducibles analyzed via minimal-degree elements in the ideal $(p, f)$, division with remainder in $Q[x]$ over the field of fractions $Q = \operatorname{Frac}(R)$, and denominator clearing.
        \end{itemize}
    \end{itemize}
\end{itemize}
\end{overview}

\section{Gauss's Theorem and Preliminaries on Content}
% Old section name: Gauss's Theorem on Polynomial Rings over UFDs

\begin{speech}[00:00:00 - 00:00:41]
Okay, so, we start, and the goal today is this Theorem B that I said, which is that $R$ is a UFD implies the polynomial ring is a UFD. \inlinemetanote{writes: Thm B: $R$ is a UFD $\implies R[x]$ is a UFD} And this theorem is, generally, well, it's a theorem of Carl Friedrich Gauss \inlinemetanote{writes: (C.F. Gauss)}, whose ghost hovers over mathematics.

And it has also an argument which has some intricacy, like the last argument we had.

Doesn't it seem like there's an echo? Or not? No? Okay, good.
\end{speech}

\recap{Theorem~B was stated in the previous lecture (\cref{thm:ufd-polynomial-ring}), together with its consequence that $K[x_1, \dots, x_n]$ is a UFD (\cref{ex:multivariate-poly-ufd}). The lecturer restates it.}

\begin{content}
\begin{theorem}[Gauss's Theorem on Polynomial Rings (Theorem B)]\label[theorem]{thm:ufd-polynomial-ring-w6wed}
Let $R$ be a unique factorization domain (UFD). Then the polynomial ring in one variable $R[x]$ is also a unique factorization domain.
\end{theorem}

\begin{steps}[Overview and Significance of Gauss's Theorem]
Theorem B generalizes the classical result that if $F$ is a field, then $F[x]$ is a UFD (in fact, a Principal Ideal Domain, \cref{prop:kx-ideals-principal}). However, when $R$ is merely a UFD (such as $R = \mathbb{Z}$), the polynomial ring $R[x]$ is a PID only if $R$ is a field (for instance, the ideal $(2, x) \subset \mathbb{Z}[x]$ cannot be generated by a single element, \cref{exer:prove-2-x-not-principal}; for a non-unit $a \neq 0$ of $R$, the ideal $(a, x)$ plays the same role). Consequently, proving that $R[x]$ is nonetheless a UFD requires a delicate analysis of polynomial factorization using the concepts of \emph{content} and \emph{primitive polynomials}.
\end{steps}
\end{content}

\subsection{Irreducible Elements and Prime Ideals in a UFD}

\begin{speech}[00:00:41 - 00:01:51]
So, the first thing that I will say is an exercise which I really encourage you to do. It says that if $R$ is a UFD, then we have an equivalence --- because it's part of this language that the equivalence is that $p$ is irreducible (that is the notion we had defined) if and only if the ideal is prime. \inlinemetanote{writes: [Ex. If $R$ is a UFD, then $p$ is irred. $\iff (p)$ is prime]}

So, this is an exercise; I don't do it here. I very much encourage you to do it because it has the ideas that we used a lot in the last lecture. You have to use, in one way, the property of being a UFD, and in the other way, the definition of a prime ideal.

And this exercise allows us to conflate, if we're in a UFD, then these two notions of being irreducible and $p$ prime are the same notion. It's like the integers; the integers have this property. So, I encourage you to do that. Actually, I strongly encourage you to do this exercise. It's part of understanding the definitions, etc., in the previous lecture.
\end{speech}

\begin{content}
\begin{exercise}[Equivalence of Irreducible Elements and Prime Elements in a UFD]\label[exercise]{exer:irred-prime-ufd}
Let $R$ be a unique factorization domain, and let $p \in R \setminus \{0\}$. Prove that $p$ is an irreducible element if and only if the principal ideal $(p) = pR$ is a prime ideal (or equivalently, $p$ is a prime element).
\end{exercise}

\begin{steps}[Proof Outline for the Exercise]
In any integral domain $R$, a prime element is always irreducible:
\begin{itemize}
    \item \textbf{Prime implies Irreducible:} If $(p)$ is prime and $p = ab$, then $ab \in (p)$, so $p \mid a$ or $p \mid b$. If $p \mid a$, write $a = pr$. Then $p = prb \implies 1 = rb$, so $b$ is a unit in $R$. Thus, $p$ is irreducible.
    \item \textbf{Irreducible implies Prime (requires UFD):} Suppose $p$ is irreducible and $p \mid ab$. Then $ab = pc$ for some $c \in R$. If $a = 0$ or $b = 0$, there is nothing to prove; if $a$ is a unit, then $p \mid a^{-1} a b = b$, and similarly if $b$ is a unit. Otherwise, factoring $a, b$ (and $c$, unless it is a unit) into products of irreducible elements, the irreducible element $p$ appearing on the right-hand side must, by uniqueness of factorization up to associates, be associate to one of the irreducible factors of $a$ or of $b$. Hence, $p \mid a$ or $p \mid b$, which implies $(p)$ is a prime ideal.
\end{itemize}
The second half is the generalization of Euclid's lemma (\cref{fact:pre-euclid}) to UFDs; for PIDs it was proved directly as Step~2 of Theorem~A (\cref{lem:step2-pid-irr-implies-prime}). The first half, together with Step~1 of Theorem~A (\cref{clm:step1-uniqueness}), shows: a domain in which every non-zero non-unit has a factorization into irreducibles is a UFD if and only if every irreducible element is prime.
\end{steps}
\end{content}

\subsection{The Content of a Polynomial}

\begin{speech}[00:01:51 - 00:02:55]
Okay, so, there's a new notion here which is content, which is an extremely useful notion. \inlinemetanote{writes: Content} And the idea here is if I have a polynomial here --- and for this discussion of content, I don't consider the zero polynomial \inlinemetanote{writes: $f \in R[x], f \neq 0$} ---, I write $f$ as $a_0$, a constant term, a linear term, and so on, and at some point it stops, because a polynomial always stops. \inlinemetanote{writes: $f = a_0 + a_1 x + \dots + a_n x^n$}

Then, this is the definition: the content --- so, the word for that is the content --- is equal to the $\operatorname{GCD}$ of these elements here: $a_0, a_1, \dots, a_n$. \inlinemetanote{writes: $\operatorname{Def.} c(f) = \operatorname{GCD}(a_0, a_1, \dots, a_n)$} Okay?

So, why are we allowed to consider $\operatorname{GCD}$? It's because this $R$ is a UFD. I should say this content, and for the whole lecture today, $R$ is a UFD. \inlinemetanote{writes: $R$ is a UFD above} Okay?
\end{speech}

\begin{content}
\begin{definition}[Content of a Polynomial]\label[definition]{def:content-poly}
Let $R$ be a unique factorization domain, and let $f \in R[x] \setminus \{0\}$ be a non-zero polynomial written as:
\[
f(x) = a_0 + a_1 x + a_2 x^2 + \dots + a_n x^n \in R[x].
\]
The \newterm{content} of $f$, denoted $c(f)$, is defined as the greatest common divisor of all the non-zero coefficients of $f$:
\begin{equation}\label{eq:content-def}
c(f) := \operatorname{GCD}(a_0, a_1, \dots, a_n).
\end{equation}
\end{definition}
\end{content}

\begin{speech}[00:02:55 - 00:04:53]
And then, this is how it's typically written, and what do I mean by this? Well, we already discussed what $\operatorname{GCD}$ is. You're not allowed to have any zeros here. So, if any of these coefficients are zero, you just strike them out, because we don't do $\operatorname{GCD}$ of zero.

So, this is how it's typically written, but what it means is you take $f$, you put here the non-zero coefficients \inlinemetanote{writes: non-zero under $a_0, a_1, \dots, a_n$} --- non-zero coefficients ---, and then take their $\operatorname{GCD}$s, as we discussed yesterday (Recall: \cref{def:gcd-ufd}, \cref{thm:gcd-formula-ufd}).

The way you take $\operatorname{GCD}$s is for every one of these, like this one, if it's non-zero... There has to be some of these that are non-zero because $f$ is not zero. So, there's at least one that's not zero. You put the list of all the ones that are non-zero, and for each one you take the prime factorization, which is, you know, some unit and then some product of $p_i$ to the $k_i$ --- let's say $k_{i0}$ for the zero one \inlinemetanote{writes: $a_0 = u_0 \prod p_i^{k_{i0}}$}, and then $a_1$ would be some unit one, $p_i$ to the $k_{i1}$ \inlinemetanote{writes: $a_1 = u_1 \prod p_i^{k_{i1}}$}.

And then to get the $\operatorname{GCD}$ --- it's the greatest common divisor ---, you take the minimal power in each prime. We discussed it for two, but you can also discuss it for several. It's just like for the integers.

Okay? So, that's the definition of the content: you just take your polynomial, put the coefficients which are not zero, take the $\operatorname{GCD}$ of them, which is well-defined since it's a UFD.

And there's another unpleasant ambiguity here: these $\operatorname{GCD}$s are always defined only up to unit. So, this thing is defined up to unit. \inlinemetanote{boxes: up to unit} That's just how it is, but people will just talk about the content, even though implicitly it's defined only up to unit.
\end{speech}

\begin{content}[Explicit Computation of Content in a UFD]
Let $R$ be a UFD. Every non-zero coefficient $a_j \neq 0$ admits a unique prime factorization:
\[
a_j = u_j \prod_{i=1}^m p_i^{k_{ij}},
\]
where $u_j \in R^\times$ is a unit, $p_1, \dots, p_m$ are pairwise non-associate irreducible (prime) elements, and $k_{ij} \ge 0$.
The greatest common divisor of several elements is defined as for two (\cref{def:gcd-ufd}: a common divisor that every common divisor divides), and the argument of \cref{thm:gcd-formula-ufd} gives it by minimal exponents. So the greatest common divisor of the non-zero coefficients is:
\begin{equation}\label{eq:gcd-prime-powers}
\operatorname{GCD}(a_0, a_1, \dots, a_n) = \prod_{i=1}^m p_i^{\min \{ k_{ij} \mid a_j \neq 0 \}}.
\end{equation}

\begin{remark}[Ambiguity Up to Units]\label[remark]{rem:content-unit-ambiguity}
Because the greatest common divisor in an arbitrary UFD is uniquely determined only up to multiplication by an invertible element (unit) in $R^\times$, the content $c(f)$ is formally defined only up to associates (\cref{def:associate-elements}) in $R$: it is well defined as an associate class, an element of the quotient set of $R$ by the equivalence relation $\sim$ (\cref{def:pre-quotient-set}), sometimes written
\[
c(f) \in R / R^\times.
\]
By standard convention, we write $c(f) = d$ to mean that $c(f)$ is associate to $d$, or we fix a canonical choice of associates whenever one is naturally available (such as positive integers in $\mathbb{Z}$ or monic polynomials in $F[x]$).
\end{remark}
\end{content}

\subsection{Canonical Decomposition of Polynomials}

\begin{speech}[00:04:53 - 00:06:10]
All right, some properties of the content. If I take this $f$ --- and whenever I consider contents, I always have $f \neq 0$ ---, then I can write $f$ as the content times $f^*$ \inlinemetanote{writes: $f \neq 0, f = c(f) \cdot f^*$}, because since it's a $\operatorname{GCD}$, I can divide all the coefficients by the content, and I can write this equation. So, $f$ is in polynomials in $R$, this content is a number, so it's in $R$, and this is also in polynomials in $R$ \inlinemetanote{writes: $R[x]$ under $f$, $R$ under $c(f)$, and $R[x]$ under $f^*$}. Since the content is the $\operatorname{GCD}$, I can just divide it out of all of them, and I get a polynomial here. Okay?

And if you think about this, what happens for this polynomial? What is the content of this polynomial? What do you think the answer to that is? It's just like I take the coefficients, they have a $\operatorname{GCD}$, I divide out the $\operatorname{GCD}$... yeah, one, yeah. \inlinemetanote{writes: What is $c(f^*) = 1$} So, the answer is, this is one. And whenever I say something like this or whenever it's written, people will just write this, but what is meant is that this is up to multiplication by a unit, one, because that depends on how you did it. Because you've just taken out all of the common divisors, what's left is just units.
\end{speech}

\begin{content}
\begin{proposition}[Factorization into Content and Primitive Part]\label[proposition]{prop:canonical-decomp}
Let $R$ be a UFD and let $f \in R[x] \setminus \{0\}$. Then $f$ can be factored as:
\begin{equation}\label{eq:f-cf-fstar}
f = c(f) \cdot f^*,
\end{equation}
where $c(f) \in R$ and $f^* \in R[x]$. Furthermore, the content of $f^*$ satisfies:
\begin{equation}\label{eq:c-fstar-one}
c(f^*) \sim 1 \quad (\text{up to a unit in } R).
\end{equation}
\end{proposition}

\begin{steps}[Algebraic Justification]
Let $f(x) = \sum_{j=0}^n a_j x^j$ with $c(f) = d \in R$. By definition of the greatest common divisor, $d \mid a_j$ for all $j \in \{0, \dots, n\}$. Therefore, each quotient $a_j^* := a_j / d$ belongs to $R$, defining the polynomial:
\[
f^*(x) := \sum_{j=0}^n a_j^* x^j \in R[x].
\]
If an irreducible element $p \in R$ divided every non-zero coefficient $a_j^*$ of $f^*$, then $p \cdot d$ would be a common divisor of all $a_j = d \cdot a_j^*$ in $R$, hence would divide $d$: $d = p d r$, and cancelling $d \neq 0$ gives $1 = p r$, so $p$ would be a unit, which an irreducible is not. Thus no irreducible divides all $a_j^*$; since every non-unit has an irreducible factor, $\operatorname{GCD}(a_0^*, \dots, a_n^*) \sim 1$.
\end{steps}
\end{content}

\begin{speech}[00:06:10 - 00:07:13]
So, example. Let's say $R$ is the integers, that's maybe our favorite ring that's a UFD, and I take some polynomial in the integers. Okay, let's say $17x^{17} + 289$. \inlinemetanote{writes: Example: $R = \mathbb{Z}, f = 17x^{17} + 289$} Okay, how about this polynomial? What's the content of this polynomial?

Well, you can take the $17$ out --- this is maybe excessive $17$s: $x^{17} + 17$, because $289$ is the square of $17$, as you surely know. \inlinemetanote{writes: $= 17(x^{17} + 17)$} So, the content is $17$. \inlinemetanote{writes: $c(f) = 17$}

So, it's not a complicated notion, but it's going to be important. Whenever we have a polynomial with coefficients in $R$, we can somehow take out the $\operatorname{GCD}$ of all the coefficients, just like we did here --- just take it out. That's all this content is.
\end{speech}

\begin{content}
\begin{example}[Computing Content over the Integers]\label[example]{ex:content-integers}
Let $R = \mathbb{Z}$, which is a UFD. Consider the polynomial:
\[
f(x) = 17 x^{17} + 289 \in \mathbb{Z}[x].
\]
Since $289 = 17^2$, the coefficients are $a_{17} = 17$, $a_0 = 289$, and all intermediate coefficients are $0$. Factoring out the greatest common divisor of the non-zero coefficients yields:
\[
f(x) = 17 \cdot (x^{17} + 17).
\]
Therefore, the content is $c(f) = 17$ (up to the units $\pm 1$), and the remaining factor is $f^*(x) = x^{17} + 17$.
\end{example}
\end{content}

\begin{hidden}
% timestamp: 00:07:13
% topic: Content of polynomials and canonical decomposition f = c(f) f*
% board_state: thm:ufd-polynomial-ring-w6wed, exer:irred-prime-ufd, def:content-poly, prop:canonical-decomp, ex:content-integers
% next_goal: Define primitive polynomials and state Gauss's Lemma
% open_loops: none
\end{hidden}

\subsection{Primitive Polynomials}

\begin{speech}[00:07:13 - 00:08:12]
So, there's a definition that's important. Definition: $g \in R[x]$, $g \neq 0$, is primitive --- that's the word that's being defined --- if the content of $g$ is equal to one, and of course I always mean up to a unit, okay? \inlinemetanote{writes: Def. $g \in R[x], g \neq 0$ is primitive if $c(g) = 1$ (up to unit)}

That's the notion, that's the definition of primitive. Primitive are those polynomials where, if I take their coefficients, there's no $\operatorname{GCD}$. Like this one here. \inlinemetanote{points at $x^{17} + 17$} $17$ is $\operatorname{GCD}$; this one has no $\operatorname{GCD}$, because of a one and a $17$.

So, in particular, if I take this $f \neq 0$, and I take $f = c(f) \cdot f^*$, this one is primitive. \inlinemetanote{writes: \qt{primitive} under $f^*$} Okay?
\end{speech}

\begin{content}
\begin{definition}[Primitive Polynomial]\label[definition]{def:primitive-poly}
Let $R$ be a UFD. A non-zero polynomial $g \in R[x] \setminus \{0\}$ is called \newterm{primitive} if its content is a unit in $R$:
\begin{equation}\label{eq:primitive-def}
c(g) \sim 1.
\end{equation}
In other words, the non-zero coefficients of $g$ share no common non-unit divisors in $R$.
\end{definition}

\begin{observation}[Primitive Factor in Canonical Decomposition]\label[observation]{obs:fstar-primitive}
For any non-zero polynomial $f \in R[x] \setminus \{0\}$, the polynomial factor $f^*$ in the canonical decomposition:
\[
f = c(f) \cdot f^*
\]
is always a primitive polynomial in $R[x]$. For instance, in \cref{ex:content-integers}, the polynomial $x^{17} + 17 \in \mathbb{Z}[x]$ is primitive because $\operatorname{GCD}(1, 17) = 1$.
\end{observation}
\end{content}

\begin{speech}[00:08:12 - 00:10:10]
And the next simple observation is: suppose $f = \gamma \cdot g$, where $\gamma$ is a scalar and $g$ is a polynomial that is primitive. \inlinemetanote{writes: Suppose $f = \gamma \cdot g, \gamma \in R, g \in R[x] \text{ primitive}$} Okay, suppose I have any expression, where $f$ is just some polynomial in $R[x]$, always not zero, and I write this polynomial as some scalar times a polynomial which is primitive. We already know one way to do that. Here, I've written it as some scalar times a polynomial that's primitive.

And the answer is that if I do this, it must be the case that this content is equal to $\gamma$ up to a unit, and $f^*$ is equal to $g$, again, up to a unit. \inlinemetanote{writes: $\implies c(f) \sim \gamma, f^* \sim g$} And the reason is because this is the $\operatorname{GCD}$ of all the coefficients. If I'm able to take this scalar out, then this scalar must divide the $\operatorname{GCD}$, because it's a divisor and it should divide the greatest common divisor. And if there's anything left over, then this wouldn't be primitive.

So, you just have to think about it. There's actually nothing going on in these statements except understanding this move of taking out the $\operatorname{GCD}$ from all of the coefficients. So, this is going to be useful for us.

Those are the basic facts about content: this content exists, and when you take $f$, you can take out the content, and then you get a primitive. And anytime you see $f$ as a scalar times a primitive, it's essentially the content and the primitive associated with the content. Okay. So, this is going to play a role in the proof.
\end{speech}

\begin{content}
\begin{proposition}[Uniqueness of Scalar and Primitive Factorization]\label[proposition]{prop:primitive-decomp-uniqueness}
Let $R$ be a UFD and let $f \in R[x] \setminus \{0\}$. Suppose that $f$ is written as:
\begin{equation}\label{eq:f-gamma-g}
f = \gamma \cdot g,
\end{equation}
where $\gamma \in R \setminus \{0\}$ and $g \in R[x]$ is a primitive polynomial. Then:
\begin{equation}\label{eq:gamma-is-content}
\gamma \sim c(f) \quad \text{and} \quad g \sim f^* \quad (\text{up to units in } R).
\end{equation}
\end{proposition}

\begin{proof}
Let $g(x) = \sum_{j=0}^n b_j x^j$. Then $f(x) = \sum_{j=0}^n (\gamma b_j) x^j$. Computing the content of $f$ using the properties of the greatest common divisor yields:
\[
c(f) = \operatorname{GCD}(\gamma b_0, \dots, \gamma b_n) \sim \gamma \cdot \operatorname{GCD}(b_0, \dots, b_n) = \gamma \cdot c(g).
\]
Since $g$ is primitive, $c(g) \sim 1$, which immediately implies:
\[
c(f) \sim \gamma \cdot 1 = \gamma.
\]
Writing $\gamma = u \cdot c(f)$ for some unit $u \in R^\times$, substituting back into \eqref{eq:f-gamma-g} gives:
\[
c(f) \cdot f^* = f = u c(f) \cdot g \implies c(f)(f^* - u g) = 0.
\]
Since $R[x]$ is an integral domain and $c(f) \neq 0$, it follows that $f^* = u g$, so $g \sim f^*$.
\end{proof}
\end{content}

\subsection{Gauss's Lemma}

\begin{speech}[00:10:10 - 00:11:20]
So, then there's two observations. The first observation is by Gauss, and this is a rather famous observation. It says that suppose $f$ and $g$ are primitive, then $f \times g$ is primitive. \inlinemetanote{writes: $\implies$ Gauss: Suppose $f, g \in R[x]$ are primitive, then $fg \in R[x]$ is primitive}

Remember, primitive means that the coefficients of the polynomial have no non-trivial $\operatorname{GCD}$; the $\operatorname{GCD}$ is just a unit. And it says if I take two of those and I multiply them, it stays primitive.

So, this is saying something extremely explicit. It's saying, for example, for integers, if you have a polynomial with integer coefficients --- two of them --- such that this one has no $\operatorname{GCD}$ --- well, the $\operatorname{GCD}$ of its coefficients is one --- and that one has a $\operatorname{GCD}$ of coefficients one, when you multiply them, the resulting polynomial has $\operatorname{GCD}$ of its coefficients one. That's something extremely explicit.

Okay, the proof of this by now for us is also very simple.
\end{speech}

\begin{content}
\begin{lemma}[Gauss's Lemma]\label[lemma]{lem:gauss-lemma}
Let $R$ be a unique factorization domain. If $f, g \in R[x]$ are primitive polynomials, then their product $fg \in R[x]$ is also a primitive polynomial:
\begin{equation}\label{eq:gauss-lemma}
c(f) \sim 1 \quad \text{and} \quad c(g) \sim 1 \implies c(fg) \sim 1.
\end{equation}
\end{lemma}
\end{content}

\begin{speech}[00:11:20 - 00:13:25]
So, we're going to proceed by contradiction. So, assume that the content of the product is not a unit. \inlinemetanote{writes: Proof: Contradiction. Assume $c(fg)$ is not a unit.} All right?

Then there exists some prime that divides this content by the unique factorization property. \inlinemetanote{writes: Then $\exists p \in R \text{ prime s.t. } p \mid c(fg) \text{ in UFD}$} Because if I have some element that's not a unit, I can factorize it. Primitive stuff is always non-zero, because to be primitive you're not zero. So, if it were not primitive, then this would not be a unit, because being primitive is if and only if this is a unit. So, then there exists some prime that divides it.

And then comes the idea in the proof: for this prime is irreducible, so that implies that the ideal is a prime ideal --- that's that exercise (Recall: \cref{exer:irred-prime-ufd}). \inlinemetanote{writes: $p \text{ is irred} \iff (p) \text{ is prime ideal}$} And then we do something interesting, which is we look at this map, a homomorphism from $R[x]$ to... So, there's a quotient map $R$ to $R \pmod p$. \inlinemetanote{writes: $R \to R/(p)$} That's the quotient map from previous weeks (Recall: \cref{def:canonical-quotient-map}). That's a homomorphism. And that also gives me a homomorphism from polynomials by just applying this quotient map to the coefficients of the polynomials. \inlinemetanote{writes: $R[x] \to (R/(p))[x]$} And this is a ring homomorphism.

Now, this is an integral domain. \inlinemetanote{writes: int. domain under $R/(p)$ and $(R/(p))[x]$} Why is it an integral domain? It's because this is a prime ideal. Prime ideal means this is an integral domain (Recall: \cref{prop:quotient-integral-domain-prime}). All right?
\end{speech}

\begin{speech}[00:13:25 - 00:15:16]
So, now I consider $\phi(f) \cdot \phi(g)$. That's equal to, because of the homomorphism property, $\phi(fg)$. \inlinemetanote{writes: $\phi(f) \cdot \phi(g) = \phi(fg)$} And by this hypothesis that $p$ divides all the coefficients of $fg$, this is equal to zero. \inlinemetanote{writes: $= 0$ (since $p \mid \text{all coeffs})$} Since $p$ divides all the coefficients of $fg$, it's equal to zero.

That means, since this is an integral domain... If I take a polynomial ring in one variable over an integral domain, the whole ring is an integral domain (Recall: \cref{thm:domain-poly-power-series}). So, we have two things multiplying to zero in an integral domain; that means either $\phi(f) = 0$ or --- let me erase that --- $\phi(g) = 0$. \inlinemetanote{writes: $\implies \phi(f) = 0 \text{ or } \phi(g) = 0$}

One of these two holds, but that contradicts the assumption. Because for either of these to hold, that means $p$ divides all their coefficients. But the assumption is that these are primitive, so no irreducible divides all their coefficients.

So, that's the proof. It's a simple proof. Gauss must have been very happy when he thought of it. And it's something, yeah, you can think about.

To understand this proof, you have to remember all these constructions: that given this ring, I take this prime ideal, I get an integral domain; integral domains stay integral when I take a polynomial ring, and then I use that property here.
\end{speech}

\begin{content}
\begin{proof}[Proof of \cref{lem:gauss-lemma}]
We proceed by contradiction. Suppose that $fg$ is not primitive. Then its content $c(fg)$ is not a unit in $R$.
Since $R$ is a UFD and $c(fg) \neq 0$ is a non-unit, there exists an irreducible element $p \in R$ such that:
\[
p \mid c(fg).
\]
By \cref{exer:irred-prime-ufd}, because $R$ is a UFD, the irreducible element $p$ generates a prime ideal $(p) \subset R$. Consequently, the quotient ring:
\[
\overline{R} := R / (p)
\]
is an integral domain (\cref{prop:quotient-integral-domain-prime}).
The canonical projection homomorphism $\pi \colon R \to R / (p)$ (\cref{def:canonical-quotient-map}) induces a natural ring homomorphism on polynomial rings (by the universal property, \cref{thm:universal-property-polynomial-ring}, sending $x \mapsto x$; it reduces every coefficient):
\begin{align*}
\phi \colon R[x] &\to (R/(p))[x], \\
\sum_{j=0}^m c_j x^j &\mapsto \sum_{j=0}^m \overline{c_j} x^j.
\end{align*}
Since $\overline{R}$ is an integral domain, the polynomial ring $\overline{R}[x] = (R/(p))[x]$ is likewise an integral domain (\cref{thm:domain-poly-power-series}).
Applying $\phi$ to the product $fg$ and using the ring homomorphism property yields:
\[
\phi(f) \cdot \phi(g) = \phi(fg).
\]
Because $p \mid c(fg)$, the prime $p$ divides every coefficient of $fg$. Therefore, every coefficient of $fg$ maps to $0$ in $R/(p)$, which implies:
\[
\phi(fg) = 0 \in (R/(p))[x].
\]
Thus, in the integral domain $(R/(p))[x]$, we have:
\[
\phi(f) \cdot \phi(g) = 0 \implies \phi(f) = 0 \quad \text{or} \quad \phi(g) = 0.
\]
If $\phi(f) = 0$, then every coefficient of $f$ is divisible by $p$, meaning $p \mid c(f)$, which contradicts the assumption that $f$ is primitive ($c(f) \sim 1$). Similarly, if $\phi(g) = 0$, then $p \mid c(g)$, contradicting the assumption that $g$ is primitive.
This contradiction proves that $fg$ must be primitive.
\end{proof}
\end{content}

\begin{hidden}
% timestamp: 00:15:16
% topic: Gauss's Lemma proven via reduction modulo a prime ideal (p)
% board_state: thm:ufd-polynomial-ring-w6wed, lem:gauss-lemma, def:primitive-poly, prop:primitive-decomp-uniqueness
% next_goal: Prove divisibility property of primitive polynomials
% open_loops: none
\end{hidden}

\subsection{A Divisibility Property of Primitive Polynomials}

\begin{speech}[00:15:16 - 00:16:14]
\inlinemetanote{The lecturer wipes the right-hand boards clean with a sponge.}
Okay, and there's one more property of contents that we have to discuss before I go to the proof of Gauss's theorem, Theorem B. Oh, it's already gone.

So, the first property was Gauss; the second property that I need is the following claim, which is also a kind of pretty claim. \inlinemetanote{writes: Claim: $g$ is primitive, $g \in R[x]$} Suppose that $g$ is primitive, so $g$ is in here, and $g$ divides $b \times f$ \inlinemetanote{writes: and $g \mid bf$}, where $b$ is a scalar that's not equal to zero \inlinemetanote{writes: $b \in R, b \neq 0$}, and $f$ is some element here \inlinemetanote{writes: $f \in R[x]$}. So, these are the assumptions: it's primitive, and this assumption.

Then the conclusion is that $g$ divides $f$. \inlinemetanote{writes: then $g \mid f$} Okay?
\end{speech}

\begin{content}
\begin{claim}[Divisibility Property of Primitive Polynomials]\label[claim]{clm:primitive-divisibility}
Let $R$ be a UFD. Suppose $g \in R[x]$ is a primitive polynomial, and suppose that:
\begin{equation}\label{eq:claim-hyp}
g \mid bf \quad \text{in } R[x],
\end{equation}
for some non-zero scalar $b \in R \setminus \{0\}$ and some polynomial $f \in R[x]$. Then:
\begin{equation}\label{eq:claim-concl}
g \mid f \quad \text{in } R[x].
\end{equation}
\end{claim}
\end{content}

\begin{speech}[00:16:14 - 00:17:55]
The first time you see these things, you just have to follow, but then you can think about it and understand the arguments for this. So, let us try to prove this.

We know that $g$ divides $bf$. That's the assumption. This is just language: $g$ divides $bf$ means that there exists some $h$ such that $bf = gh$ --- maybe $hg$, doesn't matter, but okay. \inlinemetanote{writes: Proof: $g \mid bf \implies \exists h \in R[x]: bf = hg$} That's what it means to say that $g$ divides $bf$. All right.

And now, we take the content of both sides. Well, take the content of this side: it says this polynomial here is the content of $bf$ times $f^*$, and I can take $h$ and take the content of $h$ times $h^* g$. \inlinemetanote{writes: $c(bf) f^* = c(h) h^* g$} Okay? So, this is just rewriting the left-hand side and the right-hand side in the contents.

Now, by the lemma of Gauss, we think about this object here: $h^* g$. $h^*$ is primitive because that's what you do when you take $h$. $h$ may not be primitive, but we took the content of $h$ out, so $h^*$ is primitive. $g$ is primitive by assumption. So, then by Gauss upstairs (Recall: \cref{lem:gauss-lemma}), $h^* \times g$ is primitive. \inlinemetanote{writes: $h^* g: h^* \text{ primitive}, g \text{ primitive} \implies h^* g \text{ is primitive by Gauss}$} Okay?
\end{speech}

\begin{speech}[00:17:55 - 00:19:24]
So, then we look at this equation: this is primitive, and this is primitive. And I said that whenever we have such an equation where we have the primitive part times the content, then the only possibility is that the primitives are equal up to a unit (Recall: \cref{prop:primitive-decomp-uniqueness}). This is a content equation: this is a primitive times some content. That was one of the comments up there.

So, this means that $f^*$, up to units, is this. \inlinemetanote{writes: $f^* \sim h^* g$} The consequence is that $g$ divides $f^*$, because that's what this equation means. And this means $g$ divides $f$, since $f = c(f) f^*$ (Recall: \cref{prop:canonical-decomp}). \inlinemetanote{writes: so $g \mid f^* \implies g \mid f \text{ since } f = c(f) f^*$} If you divide $f^*$, you also divide $f$.

So, that's the proof of this too. Now with these preliminaries, we can finally prove Gauss's theorem.
\end{speech}

\begin{content}
\begin{proof}[Proof of \cref{clm:primitive-divisibility}]
By assumption, $g \mid bf$ in $R[x]$, which means there exists a polynomial $h \in R[x]$ such that:
\begin{equation}\label{eq:bf-hg}
bf = hg.
\end{equation}
We express each side in terms of its canonical decomposition into content and primitive part:
\begin{align*}
\text{Left-hand side:} \quad bf &= c(bf) \cdot (bf)^* \sim b \cdot c(f) \cdot f^*, \\
\text{Right-hand side:} \quad hg &= (c(h) \cdot h^*) \cdot g = c(h) \cdot (h^* g).
\end{align*}
By \cref{obs:fstar-primitive}, $h^*$ is primitive. Since $g$ is primitive by assumption, Gauss's Lemma (\cref{lem:gauss-lemma}) implies that the product $h^* g$ is primitive:
\[
c(h^* g) \sim 1.
\]
Equating the two decompositions gives:
\[
c(bf) \cdot f^* = c(h) \cdot (h^* g).
\]
On both sides, we have a scalar in $R$ multiplied by a primitive polynomial ($f^*$ on the left, and $h^* g$ on the right). By the uniqueness of the primitive decomposition (\cref{prop:primitive-decomp-uniqueness}), the primitive parts must be associate in $R[x]$:
\[
f^* \sim h^* g.
\]
Hence, $g \mid f^*$ in $R[x]$. Since $f = c(f) \cdot f^*$, it follows immediately that:
\[
g \mid f \quad \text{in } R[x].
\]
\end{proof}
\end{content}

\begin{speech}[00:19:24 - 00:20:34]
It's a new idea, this content, and I've given you the definition and various properties of it. But it takes a little bit of time to understand how to use them and how to think about them, and you have to spend that time. But the idea of the content is in some sense a very simple idea: if you have this polynomial with coefficients in a UFD, you just take out the common divisor of all the coefficients, and you put it out.

And these properties I'm talking about are something about what happens to what's left. These primitive polynomials become a little bit special.

So, now we come to Theorem B, which is the second theorem for the week. And it also, as I said, has a little bit of an intricate logic. I did look at the book to try to follow more or less the notation. It's not going to be perfect, but I did try to at least match the notation. \inlinemetanote{The lecturer finishes erasing the upper chalkboard.}
\end{speech}

\begin{hidden}
% timestamp: 00:20:34
% topic: Preliminaries complete; transition to Proof of Theorem B (R UFD => R[x] UFD)
% board_state: thm:ufd-polynomial-ring-w6wed, clm:primitive-divisibility, lem:gauss-lemma
% next_goal: State Theorem B and prove Step 1 (Existence of factorization into irreducibles)
% open_loops: none
\end{hidden}

\section{Proof of Gauss's Theorem}

\begin{speech}[00:20:34 - 00:21:44]
Okay, so, we want to prove: $R$ is a UFD implies $R[x]$ is a UFD. \inlinemetanote{writes: Thm. $R \text{ UFD} \implies R[x] \text{ is a UFD}$} So, for $R[x]$, we must prove two things: we must prove the existence statement --- existence of factorization into irreducibles in $R[x]$ \inlinemetanote{writes: for $R[x]$, we must prove: $\exists \text{ factorization into irreducibles in } R[x]$} ---, and then we must prove the uniqueness statement: uniqueness \inlinemetanote{writes: $! \text{ uniqueness}$}.

There's no shortcut to that: we have to prove existence and uniqueness to prove it's a UFD.

So, let us start this time with existence. Step 1 is the existence result. \inlinemetanote{writes: STEP 1 $\exists$}
\end{speech}

\begin{content}[Statement of Theorem B and Two-Part Strategy]
The lecturer writes Theorem~B on the board once more (\cref{thm:ufd-polynomial-ring-w6wed}):
\[
R \text{ is a UFD} \implies R[x] \text{ is a UFD}.
\]
To establish that $R[x]$ is a UFD (\cref{def:ufd}), we must prove two fundamental assertions for every non-zero, non-unit polynomial $f \in R[x]$:
\begin{enumerate}[label=\textbf{(\arabic*)}]
    \item \textbf{Existence of Factorization:} Every non-zero, non-unit $f \in R[x]$ can be factored into a product of irreducible elements in $R[x].$
    \item \textbf{Uniqueness of Factorization:} Such factorization into irreducibles is unique up to permutation and associate elements.
\end{enumerate}
\end{content}

\begin{proofWrapper}[Proof of Gauss's Theorem (\cref{thm:ufd-polynomial-ring-w6wed})]

\subsection{Step 1: Existence of Factorization}

\begin{speech}[00:21:44 - 00:23:58]
For this there's a ground case and an induction.

The ground case is: suppose $f \in R[x]$... \inlinemetanote{writes: Ground case:} Here $f$ is not zero, not a unit, and suppose the degree of $f$ is zero. \inlinemetanote{writes: Suppose $f \in R[x], f \neq 0, \text{not unit, and suppose } \deg(f) = 0$}

We're going to prove existence based on induction of the degree of the polynomial. All right? The polynomials have degree in any ring, so we're going to start with the case where the degree is zero. What does that mean, the degree is zero? If $f$ is a degree-zero polynomial, then it's like $a_0$ plus nothing else. So, this means $f$ is a scalar, actually. This means $f$ is just a scalar and it lives in $R$. \inlinemetanote{writes: $f = a_0 \implies \text{This means } f \in R$} A degree-zero polynomial is just a scalar.

So, you can imagine, in this ground case, we need to factor something in $R$. How are we going to do that? Well, we can't do that; we just assumed that from the beginning: $R$ is a unique factorization domain. So, we can just factor it.

We can factor $f$, since it's a scalar, into some irreducible factors: \inlinemetanote{writes: We can factor $f$, since it's a scalar, $f = q_1 \dots q_m$, irred. in $R$} irreducibles in $R$. And every irreducible in $R$ --- these are also irreducibles in $R[x]$, \inlinemetanote{writes: these are also irreducibles in $R[x]$} because irreducibility is a condition on what happens when you factor. And if you have a scalar, when you factor you don't get any $x$'s, because the $x$'s would never go.

All right? So, that's the ground case. In the ground case, we have degree zero, and then it's just a scalar. We don't really do anything; all we do is just take this assumption: it's a UFD, so we can factor the scalars. So, actually nothing happened in this ground case.
\end{speech}

\begin{content}[{Existence: Ground Case \texorpdfstring{($\deg(f) = 0$)}{(deg(f) = 0)}}]
\textbf{Ground Case ($\deg(f) = 0$):}
Suppose $f \in R[x]$ is a non-zero, non-unit polynomial with $\deg(f) = 0$.
Then $f = a_0 \in R$. Since $f$ is not a unit in $R[x]$, $a_0$ is a non-zero, non-unit element of $R$.

Because $R$ is a UFD by assumption, $a_0$ factors into a product of irreducible elements in $R$:
\[
f = q_1 \cdot q_2 \cdots q_m, \quad \text{where each } q_i \in R \text{ is irreducible in } R.
\]
Furthermore, every element $q \in R$ that is irreducible in $R$ remains irreducible as a constant polynomial in $R[x]$.

\begin{steps}[{Why Irreducibles in $R$ Remain Irreducible in $R[x]$}]
Suppose an irreducible element $q \in R$ factors in $R[x]$ as $q = g(x) h(x)$ for some $g, h \in R[x]$.
Since $R$ is an integral domain, the degree of a product equals the sum of degrees (the product of the leading coefficients is non-zero):
\[
0 = \deg(q) = \deg(g) + \deg(h) \implies \deg(g) = 0 \quad \text{and} \quad \deg(h) = 0.
\]
Thus, $g, h \in R$. Because $q$ is irreducible in $R$, either $g$ or $h$ must be a unit in $R$. Since the units of $R[x]$ are precisely the units of $R$ (\cref{thm:units-polynomial-ring-integral-domain}), it follows that either $g$ or $h$ is a unit in $R[x]$. Hence, $q$ remains irreducible as an element of $R[x]$.
\end{steps}
\end{content}

\begin{speech}[00:23:58 - 00:25:57]
The induction case is slightly more interesting.

In the induction case, suppose we have this $f \in R[x]$, not zero, not equal to zero or a unit. \inlinemetanote{writes: Induction case: Suppose $f \in R[x], f \neq 0 \text{ unit}, \deg(f) > 0$} And now the degree of $f$ is positive. We want to show that there exists a factorization into irreducibles in the polynomial ring. And now we actually have to do something. We didn't here, because that was just assumption.

Let us see.

As far as Theorem B is concerned, this is the first moment where an idea is needed. We need to prove existence. We're not necessarily a scalar, so we can't just use the assumption; we're some actual bona fide polynomial, so we have here $f$. And the idea here is we know one thing that I taught at the beginning of the class: we can write $f$ as $c(f)$ times $f^*$ (Recall: \cref{prop:canonical-decomp}). \inlinemetanote{writes: $f = c(f) \cdot f^*$} That's the content, and this is primitive. \inlinemetanote{writes: \qt{primitive} above $f^*$} Okay?

So, we need to factor $f$. We need to show there exists a factorization of $f$ into irreducibles for $R[x]$. We already have a kind of factorization here. And remember, this is a scalar \inlinemetanote{writes: $R$ under $c(f)$} and this is in $R[x]$ \inlinemetanote{writes: $R[x]$ under $f^*$}. And since the degree of $f$ is positive, the degree of this is positive. So, I just write that explicitly: the degree of $f^*$ is equal to the degree of $f$, which is positive \inlinemetanote{writes: $\deg(f^*) = \deg(f) > 0$}, because this is just a scalar, okay?

Now, we have to think: what could happen here? So, what could happen? \inlinemetanote{writes: What could happen?}
\end{speech}

\begin{content}[Induction Step: Canonical Decomposition Setup]
\textbf{Induction Step ($\deg(f) > 0$):}
Suppose $f \in R[x]$ is a non-zero, non-unit polynomial with $\deg(f) > 0$.
By \cref{prop:canonical-decomp}, $f$ can be decomposed into its content and primitive part:
\[
f = c(f) \cdot f^*,
\]
where $c(f) \in R \setminus \{0\}$ and $f^* \in R[x]$ is a primitive polynomial ($c(f^*) \sim 1$).
Since $c(f)$ is a non-zero scalar, $\deg(c(f)) = 0$. By additivity of degrees in the integral domain $R[x]$:
\[
\deg(f) = \deg(c(f)) + \deg(f^*) = 0 + \deg(f^*) \implies \deg(f^*) = \deg(f) > 0.
\]
If $c(f) \in R$ is not a unit, the UFD property of $R$ guarantees that $c(f)$ factors into a product of irreducible elements in $R$, each of which remains irreducible in $R[x]$. Therefore, to obtain a complete factorization of $f$ into irreducibles in $R[x]$, it suffices to factor the primitive polynomial $f^*$.
\end{content}

\begin{speech}[00:25:57 - 00:27:36]
Well, one thing that could happen --- the first case --- is that $f^*$ in $R[x]$ is irreducible. \inlinemetanote{writes: \textbf{(i)} $f^* \in R[x] \text{ is irred.}$} Okay?

If this happens, we're in really good shape. Because suppose this is irreducible, then we look at $f = c(f) \cdot f^*$. \inlinemetanote{writes: $f = c(f) \cdot f^*$} And in this branch, this is irreducible in $R[x]$. \inlinemetanote{writes: irred. in $R[x]$ under $f^*$} Now, this is a scalar, so we can apply the UFD property of $R$ and factor that into irreducibles. \inlinemetanote{writes: UFD in $R$, factor into irreds under $c(f)$} Okay?

So, then we're done, because we have irreducibles and irreducibles, so we are done. We are done with this case. \inlinemetanote{writes: We are done with this case}

This is not really a surprise: at the very beginning, anything in $R$ we can already factor into irreducibles, so we're interested in being able to factor that into irreducibles. If I just assume it's irreducible, then I'm done.
\end{speech}

\begin{content}[{Case (i): Primitive Factor \texorpdfstring{$f^*$}{f*} is Irreducible}]
\textbf{Case (i): $f^*$ is irreducible in $R[x]$.}
If $f^*$ is already irreducible in $R[x]$, then:
\begin{itemize}
    \item \textbf{Unit content ($c(f) \in R^\times$):} If $c(f) \in R^\times$ is a unit, then $f = c(f) f^*$ is associate to $f^*$ and is therefore itself irreducible in $R[x]$.
    \item \textbf{Non-unit content ($c(f) \notin R^\times$):} If $c(f) \notin R^\times$, then $c(f)$ is a non-zero, non-unit scalar in the UFD $R$. Thus, $c(f)$ factors into a finite product of irreducible elements in $R$:
    \[
    c(f) = q_1 \cdot q_2 \cdots q_k \quad (q_i \in R \text{ irreducible in } R).
    \]
    Since each $q_i$ remains irreducible in $R[x]$, the product:
    \[
    f = q_1 \cdot q_2 \cdots q_k \cdot f^*
    \]
    is a complete factorization of $f$ into irreducible elements in $R[x]$.
\end{itemize}
In either sub-case, the existence of a factorization into irreducibles is established.
\end{content}

\begin{speech}[00:27:36 - 00:29:34]
The other case, of course, is that $f^*$ in $R[x]$ is not irreducible. \inlinemetanote{writes: \textbf{(ii)} $f^* \in R[x] \text{ is not irred.}$} Okay, that's the case where I have to worry. And then again you can imagine what I'm going to do: if it's not irreducible, that means $f^*$ factors into some $A$ and $B$, \inlinemetanote{writes: $f^* = A \cdot B$} where $A$ and $B$ are in $R[x]$ \inlinemetanote{writes: $A, B \in R[x]$}, and $A$ is not a unit and $B$ is not a unit. \inlinemetanote{writes: $A$ is not a unit, $B$ is not a unit} Because if it's not irreducible, then there's a legitimate factorization. Okay? That's the other world: either it's irreducible, or there's a legitimate factorization.

And we would like to now use induction on it. If $A$ and $B$ have degree less than $f$, we are done by induction. \inlinemetanote{writes: If $A, B$ have degree $< \deg(f)$, we are done by induction} Because we're proceeding by induction on $f$, and if these are of lower degree, then we just use the induction statement to factor them. This is how you factor integers.

So, this is almost the whole proof, but now there's one case we have to worry about. Normally, if you factor a polynomial, these things will be of degree less, but there's one case they won't. The only remaining worry to consider, let's say it that way, \inlinemetanote{writes: Only remaining worry to consider is} is if one of these, say $A$, is a scalar --- that means the degree of $A$ is zero. \inlinemetanote{writes: $A \in R, \deg(A) = 0$} Because if the degree of $A$ is zero, then the degree of $B$ is the degree of $f^*$, and we can't use induction. Is that clear? That's actually the whole point of the argument, the only interesting point of the argument here. Up to now we've been just kind of shuffling the pawns around, but this is the actual point where something happens.
\end{speech}

\begin{content}[Case (ii): Reducible Primitive Factor and the Degree Obstacle]
\textbf{Case (ii): $f^*$ is reducible in $R[x]$.}
If $f^*$ is reducible in $R[x]$, then by definition there exist non-units $A, B \in R[x]$ such that:
\begin{equation}\label{eq:fstar-AB}
f^* = A \cdot B.
\end{equation}
Because $R$ is an integral domain, the degree is additive on products of non-zero polynomials:
\[
\deg(f^*) = \deg(A) + \deg(B).
\]
\begin{itemize}
    \item \textbf{The Desired Induction Step:} If both $\deg(A) < \deg(f^*) = \deg(f)$ and $\deg(B) < \deg(f^*) = \deg(f)$, then by the induction hypothesis, both $A$ and $B$ factor into products of irreducible elements in $R[x]$. Multiplying these factorizations together would immediately produce a factorization of $f^*$, and hence of $f = c(f) f^*$.
    \item \textbf{The Potential Degree Obstacle:} The induction hypothesis fails if one of the factors, say $A$, has degree zero:
    \[
    \deg(A) = 0 \implies A \in R.
    \]
    In that scenario, $\deg(B) = \deg(f^*) = \deg(f)$, meaning that the degree has not strictly decreased.
\end{itemize}
\end{content}

\begin{speech}[00:29:34 - 00:29:46]
So, this is the case to worry about. How can we exclude that? You tell me: how do we exclude this case?
\end{speech}

\begin{student}[Student Answer]
Because $f^*$ is primitive.
\end{student}

\begin{speech}[00:29:48 - 00:30:58]
Yes! $f^*$ is primitive. There is no common factor for all of its coefficients. If this were not a unit, this would be a common factor for all the coefficients.

So, the only way this happens... Since $f^*$ is primitive, this would imply $A$ is a unit. \inlinemetanote{writes: Since $f^*$ is primitive, $A$ is a unit} But that's a contradiction, because $A$ is not a unit; it's born not a unit. \inlinemetanote{writes: Contradiction}

The logic is: if this is reducible, I get $A$ and $B$, neither of which are units. And they both have degree strictly less than $f$, unless one of them is a scalar, and that's excluded because this is primitive. This is why that definition was made in the first place for that step in the argument.

Okay, so, this finishes the existence!

Now there's uniqueness. Uniqueness is fun, maybe more fun.
\end{speech}

\begin{content}[Resolution via Primitivity and Conclusion of Step 1]
\begin{claim}[Exclusion of Degree-Zero Factors]\label[claim]{clm:exclude-deg0-factors}
In the factorization $f^* = AB$ of the primitive polynomial $f^*$, neither factor can have degree zero.
\end{claim}

\begin{proof}[Proof of \cref{clm:exclude-deg0-factors}]
Suppose for contradiction that $\deg(A) = 0$, so $A = a \in R$.
Since $A$ is not a unit in $R[x]$ by assumption, $a$ is not a unit in $R$.
Because $a$ divides $f^* = aB$ in $R[x]$, $a$ must divide every coefficient of $f^*$.
Consequently, $a$ divides the content of $f^*$:
\[
a \mid c(f^*) \quad \text{in } R.
\]
However, $f^*$ is primitive (\cref{def:primitive-poly}), which means $c(f^*) \sim 1$. Thus, $a \mid 1$ in $R$, implying that $a \in R^\times$ is a unit in $R$.
This contradicts the fact that $A$ is not a unit in $R[x]$.
\end{proof}

\begin{steps}[Conclusion of the Induction Step for Existence]
Since degree-zero non-unit factors cannot occur, we have:
\[
\deg(A) \ge 1 \quad \text{and} \quad \deg(B) \ge 1.
\]
Because $\deg(A) + \deg(B) = \deg(f^*) = \deg(f)$, it follows that:
\[
\deg(A) < \deg(f) \quad \text{and} \quad \deg(B) < \deg(f).
\]
By the induction hypothesis on the degree, both $A$ and $B$ factor into products of irreducible elements in $R[x]$. Therefore, $f^* = AB$ factors into irreducibles, and combining this with the irreducible factorization of the scalar content $c(f) \in R$ yields a complete factorization of $f = c(f) f^*$ into irreducible elements in $R[x]$.
This completes the proof of Step 1 (Existence of Factorization).
\end{steps}
\end{content}

\begin{hidden}
% timestamp: 00:31:17
% topic: Completed Step 1 (Existence) of Gauss's Theorem; transitioning to Step 2 (Uniqueness)
% board_state: thm:ufd-polynomial-ring-w6wed, prop:canonical-decomp, lem:gauss-lemma, clm:primitive-divisibility
% next_goal: Prove that irreducibles in R[x] are prime to deduce uniqueness
% open_loops: none
\end{hidden}

\begin{speech}[00:30:58 - 00:31:28]
So, you see how cunning this definition was. As I said, if you're working on this problem, the first thing you would not do is write down the definition. It comes only after you understand the argument. Then you can go back to the beginning and write down the definition for content. \inlinemetanote{The lecturer wipes the lower-left chalkboard clean with a sponge.}
\end{speech}

\subsection{Step 2: Uniqueness of Factorization}

\begin{speech}[00:31:28 - 00:33:17]
Okay, so, now we have to prove the uniqueness, and the path to proving uniqueness was already explained in the previous lecture.

Step 2: prove uniqueness. \inlinemetanote{writes: STEP 2 Prove !} We already had shown a method to prove uniqueness: we already had shown that if $p$ irreducible implies $p$ prime, then uniqueness holds. \inlinemetanote{writes: We already had shown: if $p$ irred $\implies p$ prime $\implies$ uniqueness holds} Because that was one of our steps yesterday (Recall: \cref{clm:step1-uniqueness}): if I have any ring where irreducible implies prime, then I get uniqueness.

So, I have to prove this now. That's the last thing I have to do.

You should look at the lecture from yesterday: we must prove that if $p \in R[x]$ is irreducible, \inlinemetanote{writes: We must prove: $p \in R[x]$ irred} then it's prime. The way you do that is you have to say that we have the assumption: $p$ is irreducible, $fg \in (p)$, \inlinemetanote{writes: $fg \in (p)$} $f \notin (p)$, \inlinemetanote{writes: $f \notin (p)$} and we must prove under these hypotheses that $g \in (p)$. \inlinemetanote{writes: $\implies g \in (p)$} That's the last thing.

If we prove this --- you assume it's irreducible, these things lie in $(p)$, $f$ is not in $(p)$, and then $g$ is in $(p)$ ---, then that finishes the proof of Gauss's theorem.
\end{speech}

\begin{content}[{Strategy for Uniqueness: Irreducible Implies Prime in \texorpdfstring{$R[x]$}{R[x]}}]
\textbf{Step 2 (Uniqueness):}
By Step~1 of the proof of Theorem~A (\cref{clm:step1-uniqueness}), in any integral domain where every irreducible element is a prime element, factorization into irreducibles is unique up to permutation and associate elements.

Therefore, to prove that $R[x]$ is a UFD, it suffices to prove that every irreducible element in $R[x]$ is prime:
\begin{equation}\label{eq:irred-implies-prime-goal}
p \in R[x] \text{ is irreducible} \implies p \text{ is a prime element in } R[x].
\end{equation}
In terms of principal ideals, this means that for any $f, g \in R[x]$:
\begin{equation}\label{eq:prime-ideal-condition-w6wed}
fg \in (p) \quad \text{and} \quad f \notin (p) \implies g \in (p),
\end{equation}
which is equivalent to:
\[
p \mid fg \quad \text{and} \quad p \nmid f \implies p \mid g \quad \text{in } R[x].
\]
We divide the verification of \eqref{eq:prime-ideal-condition-w6wed} into two distinct cases based on the degree of $p$:
\begin{itemize}
    \item \textbf{Case 1:} $\deg(p) = 0$ ($p$ is a constant scalar in $R$).
    \item \textbf{Case 2:} $\deg(p) > 0$ ($p$ is a polynomial of positive degree).
\end{itemize}
\end{content}

\subsubsection{Case 1: Degree Zero Irreducibles (\texorpdfstring{$\deg(p) = 0$}{deg(p) = 0})}

\begin{speech}[00:33:17 - 00:34:58]
Okay, so, this is what we have to prove. And to start, there's a case that helps separately. Case: \inlinemetanote{writes: Case} just as before, $\deg(p) = 0$. \inlinemetanote{writes: $\deg(p) = 0$}

This case means that $p$ is a scalar: $p \in R$. \inlinemetanote{writes: $p \in R$} All right.

What do we have here? We need to prove this implication. $fg$ is equal to the content of $f$ times $f^*$ times the content of $g$ times $g^*$. \inlinemetanote{writes: $fg = c(f) f^* \cdot c(g) g^*$} I just used the content equation for $f$, the content equation for $g$, and then I put them in order here: content of $f$, content of $g$, $f^*$, $g^*$. \inlinemetanote{writes: $= c(f) c(g) f^* g^*$} All right?

And the second thing says that $p$ divides this: $p$ divides the whole thing \inlinemetanote{writes: $p \mid c(f) c(g) f^* g^*$} --- content of $f$, content of $g$, times these polynomials $f^*$, $g^*$.

Now, $f^*$ is primitive because that's what happens when you extract the content. \inlinemetanote{writes: \qt{primitive} under $f^*$} This one is also primitive \inlinemetanote{writes: \qt{primitive} under $g^*$} --- that's what happens when you extract the content. So, the whole thing is primitive. \inlinemetanote{writes: \qt{primitive} under $f^* g^*$} That's by the Gauss claim (Recall: \cref{lem:gauss-lemma}): Gauss says that the whole thing is primitive. \inlinemetanote{writes: Gauss $\implies f^* g^* \text{ is primitive}$}
\end{speech}

\begin{content}[Case 1: Decomposition of the Product]
\textbf{Case 1 ($\deg(p) = 0$):}
Suppose $p \in R[x]$ is irreducible with $\deg(p) = 0$. Then $p = c \in R$.
Since $p$ is irreducible in $R[x]$, it is irreducible in $R$ (a factorization in $R$ is one in $R[x]$, and the units agree). Because $R$ is a UFD, $p$ is a prime element in $R$ (\cref{exer:irred-prime-ufd}).

Suppose $p \mid fg$ in $R[x]$ and $p \nmid f$.
Decomposing $f$ and $g$ into their content and primitive parts:
\[
f = c(f) \cdot f^* \quad \text{and} \quad g = c(g) \cdot g^*,
\]
where $f^*, g^* \in R[x]$ are primitive, we express the product $fg$ as:
\[
fg = \bigl(c(f) f^*\bigr) \cdot \bigl(c(g) g^*\bigr) = \bigl(c(f) c(g)\bigr) \cdot (f^* g^*).
\]
By Gauss's Lemma (\cref{lem:gauss-lemma}), since $f^*$ and $g^*$ are primitive, their product:
\[
f^* g^* \in R[x] \quad \text{is also primitive}.
\]
Since $p \mid fg$, the scalar $p$ divides the polynomial $\bigl(c(f) c(g)\bigr) \cdot (f^* g^*)$ in $R[x]$.
\end{content}

\begin{speech}[00:34:58 - 00:36:57]
But that's interesting. $p$ divides... This is a scalar, that's a scalar, and it divides every coefficient of this polynomial, which is primitive.

There must be one coefficient of this polynomial which $p$ does not divide. $p$ cannot divide every coefficient of this polynomial, otherwise this polynomial would not be primitive.

So, there's some coefficient, let's say $a_{17}$. That's the coefficient \inlinemetanote{writes: $a_{17}$ is the coeff.} of $f^* g^*$ \inlinemetanote{writes: of $f^* g^*$} which $p$ does not divide. \inlinemetanote{writes: which $p$ does not divide} I could put a variable there too, but $17$ in this class is kind of a variable anyway. $p$ does not divide.

This means that $p$ does not divide $a_{17}$ \inlinemetanote{writes: $p \nmid a_{17}$}, but $p$ divides all of these things together: it divides $c(f) c(g) a_{17}$. \inlinemetanote{writes: $p \mid c(f) c(g) a_{17}$} Because it must divide all of this, but it doesn't divide $a_{17}$.

These are all scalars now, and we have unique factorization. So, $p$ doesn't divide that.

$p$ doesn't divide the content of $f$ either, because if $p$ divided the content of $f$, then $f$ would be in this ideal \inlinemetanote{writes: $f \notin (p) \implies p \nmid c(f)$}, because $f = c(f) \cdot f^*$. If $p$ divided the content of $f$, $f$ would be in that ideal, but it's not.

So, it doesn't divide that, doesn't divide that. By unique factorization, $p$ must divide the content of $g$ \inlinemetanote{writes: $p \mid c(g)$}, and this implies $p$ divides the polynomial $g$ \inlinemetanote{writes: $\implies p \mid g$}, which says that $g$ is in the ideal. \inlinemetanote{writes: $\implies g \in (p)$}

So, that handles the degree-zero case, and I have five minutes, and I will do the last case, hopefully.
\end{speech}

\begin{content}[{Case 1: Conclusion that \texorpdfstring{$g \in (p)$}{g in (p)}}]
\begin{proof}[Completion of Case 1 ($\deg(p) = 0$)]
Write the primitive polynomial $f^* g^*$ as:
\[
f^* g^* = \sum_{j=0}^N d_j x^j \in R[x].
\]
Since $f^* g^*$ is primitive, $\operatorname{GCD}(d_0, \dots, d_N) \sim 1$. Therefore, the prime element $p \in R$ cannot divide all coefficients $d_j$. Thus, there exists an index $k$ (the lecturer calls it $17$, writing $a_{17}$ for $d_k$) such that:
\[
p \nmid d_k \quad \text{in } R.
\]
Because $p \mid fg$ with $p \in R$ a constant, $p$ must divide every coefficient of $fg = c(f) c(g) (f^* g^*)$ in $R$. In particular, the $k$-th coefficient gives:
\begin{equation}\label{eq:p-divides-coeff}
p \mid c(f) c(g) d_k \quad \text{in } R.
\end{equation}
Since $p$ is a prime element in the UFD $R$, \eqref{eq:p-divides-coeff} implies:
\[
p \mid c(f) \quad \text{or} \quad p \mid c(g) \quad \text{or} \quad p \mid d_k.
\]
We already know $p \nmid d_k$. Moreover, if $p \mid c(f)$, then $c(f) = p \cdot r$ for some $r \in R$, which implies:
\[
f = c(f) f^* = p \cdot (r f^*) \implies p \mid f \quad \text{in } R[x] \implies f \in (p),
\]
directly contradicting the hypothesis that $f \notin (p)$.
Consequently, we must have:
\[
p \mid c(g) \quad \text{in } R.
\]
Writing $c(g) = p \cdot s$ for some $s \in R$, we conclude:
\[
g = c(g) g^* = p \cdot (s g^*) \implies p \mid g \quad \text{in } R[x] \implies g \in (p).
\]
This proves that every irreducible constant element $p \in R$ generates a prime ideal in $R[x]$.
\end{proof}
\end{content}

\begin{hidden}
% timestamp: 00:36:57
% topic: Case 1 (deg(p) = 0) complete; starting Case 2 (deg(p) > 0) using field of fractions Q = Fraction(R)
% board_state: thm:ufd-polynomial-ring-w6wed, eq:irred-implies-prime-goal, lem:gauss-lemma, clm:primitive-divisibility
% next_goal: Construct minimal degree polynomial m in (p, f) and apply division algorithm in Q[x]
% open_loops: none
\end{hidden}

\subsubsection{Case 2: Positive Degree Irreducibles (\texorpdfstring{$\deg(p) > 0$}{deg(p) > 0})}

\begin{speech}[00:36:57 - 00:39:16]
The last case is the most interesting. Let's try to do this.

The last case is that we have those hypotheses, the same on the board, but Case: \inlinemetanote{writes: Case} now the degree of $p$ is positive. \inlinemetanote{writes: $\deg(p) > 0$}

We make an ideal: $I$ is the ideal generated by $p$ and $f$. \inlinemetanote{writes: $I = (p, f) \subset R[x]$} This is an ideal in $R[x]$.

It's just some ideal. But we have the degree function here, and this is the first idea in the proof: choose an element here which is not zero --- we're never interested in zero here --- such that $m$ has minimal degree. \inlinemetanote{writes: choose $m \in I, m \neq 0$ such that $m$ has minimal degree}

We can do that, because you can always pick things of minimal degree; you can always pick the smallest number which occurs. It's important that $m$ is not zero.

A crucial thing is we're going to show now: we will show the degree of $m$ is zero. \inlinemetanote{writes: We will show $\deg(m) = 0$} So, it's a scalar. And that has a very clever argument.

What we would like to do is use the division algorithm. But we cannot use the division algorithm. We would very much like to use the division algorithm for polynomial rings. \inlinemetanote{writes: We would like to use the div. alg.} We can't do it, because the division algorithm requires a field (Recall: \cref{thm:poly-long-division}). So, we give ourselves a field. What field? The field of fractions!

Let $Q$ be the field of fractions. \inlinemetanote{writes: Let $Q = \text{Fraction}(R)$} $R$ is a domain; we've seen that to every domain you associate a field of fractions (Recall: \cref{def:field-of-fractions}). We let $Q$ be this field of fractions: $Q = \operatorname{Frac}(R)$.

So, then in $Q[x]$ we have the division algorithm. \inlinemetanote{writes: So in $Q[x]$ we have div. alg.}
\end{speech}

\begin{content}[{Case 2: The Ideal \texorpdfstring{$(p, f)$}{(p, f)} and the Field of Fractions \texorpdfstring{$Q$}{Q}}]
\textbf{Case 2 ($\deg(p) > 0$):}
Suppose $p \in R[x]$ is irreducible with $\deg(p) > 0$. We assume $fg \in (p)$ and $f \notin (p)$, and we must show $g \in (p)$.

Consider the ideal in $R[x]$ generated by $p$ and $f$:
\begin{equation}\label{eq:ideal-I-def}
I := (p, f) = \{ A p + B f \mid A, B \in R[x] \} \subseteq R[x].
\end{equation}
Because $p \in I$ and $p \neq 0$, the set of degrees of non-zero polynomials in $I$:
\[
\{ \deg(h) \mid h \in I \setminus \{0\} \} \subset \mathbb{N}_0
\]
is non-empty. By the Well-Ordering Principle of $\mathbb{N}_0$, there exists an element of minimal degree:
\begin{equation}\label{eq:m-minimal-def}
m \in I \setminus \{0\} \quad \text{such that } \deg(m) \le \deg(h) \quad \text{for all } h \in I \setminus \{0\}.
\end{equation}
\textbf{Core Objective:} We claim that $\deg(m) = 0$, meaning $m \in R \setminus \{0\}$ is a scalar.

\begin{steps}[Why We Introduce the Field of Fractions $Q$]
In $F[x]$, where $F$ is a field, division with remainder (\cref{thm:poly-long-division}) shows that an element of minimal degree in an ideal generates the entire ideal (\cref{prop:kx-ideals-principal}). However, $R[x]$ is not a PID when $R$ is not a field.
To exploit division with remainder, we embed $R$ into its field of fractions (\cref{def:field-of-fractions}):
\[
Q := \operatorname{Frac}(R).
\]
Then $R[x] \subset Q[x]$, and $Q[x]$ has division with remainder.
\end{steps}
\end{content}

\begin{speech}[00:39:16 - 00:41:26]
What does that mean? We can write $f = \hat{q} m + \hat{r}$, \inlinemetanote{writes: $f = \hat{q} m + \hat{r}$} where these $\hat{q}, \hat{r}$ are in the field of fractions. \inlinemetanote{writes: $\hat{q}, \hat{r} \in Q[x]$} And the division algorithm has the usual properties that $\hat{r}$ is either zero or the degree of $\hat{r}$ is strictly less than the degree of $m$. \inlinemetanote{writes: $\hat{r} = 0 \text{ or } \deg(\hat{r}) < \deg(m)$}

And we can also do it with $p$: $p = \hat{\hat{q}} m + \hat{\hat{r}}$. \inlinemetanote{writes: $p = \hat{\hat{q}} m + \hat{\hat{r}}$} Same thing: $\hat{\hat{q}}, \hat{\hat{r}} \in Q[x]$. \inlinemetanote{writes: $\hat{\hat{q}}, \hat{\hat{r}} \in Q[x]$} Maybe I just call it $Q$, I don't need to put that here. And again, now with $\hat{\hat{r}} = 0$ or the degree of $\hat{\hat{r}}$ is less than the degree of $m$. \inlinemetanote{writes: $\hat{\hat{r}} = 0 \text{ or } \deg(\hat{\hat{r}}) < \deg(m)$}

So, that's one of the ideas: we use the division algorithm in the field of fractions.

But we want to be in $R$! We want to be in the ring $R$, and these things can have denominators. And it's not a problem: we just multiply by the denominators to clear the denominators. By clearing the denominators, we will multiply by some $b \in R$. \inlinemetanote{writes: $b \in R$}

So, if we just multiply by the denominators, I get an equation: $b' f = q' m + r'$, \inlinemetanote{writes: $b' f = q' m + r'$} and similarly $b'' p = q'' m + r''$. \inlinemetanote{writes: $b'' p = q'' m + r''$}

And now, if you look at this, everything is in $R$: all the terms here are in the ring $R$. \inlinemetanote{writes: $b', b'' \in R; q', q'', r', r'' \in R[x]$}
\end{speech}

\begin{content}[{Division with Remainder in \texorpdfstring{$Q[x]$}{Q[x]} and Clearing Denominators}]
Applying division with remainder in $Q[x]$ (\cref{thm:poly-long-division}, with $K = Q$) to divide $f$ and $p$ by $m \neq 0$, there exist unique quotients and remainders $\hat{q}, \hat{r}, \hat{\hat{q}}, \hat{\hat{r}} \in Q[x]$ such that:
\begin{align}
f &= \hat{q} m + \hat{r}, && \text{with } \hat{r} = 0 \text{ or } \deg(\hat{r}) < \deg(m), \label{eq:div-f-Q} \\
p &= \hat{\hat{q}} m + \hat{\hat{r}}, && \text{with } \hat{\hat{r}} = 0 \text{ or } \deg(\hat{\hat{r}}) < \deg(m). \label{eq:div-p-Q}
\end{align}
Since the coefficients of $\hat{q}, \hat{r}$ and $\hat{\hat{q}}, \hat{\hat{r}}$ are fractions in $Q = \operatorname{Frac}(R)$, we clear the denominators. By multiplying \eqref{eq:div-f-Q} by a common denominator $b' \in R \setminus \{0\}$ of all coefficients of $\hat{q}$ and $\hat{r}$ (for instance the product of their denominators) and multiplying \eqref{eq:div-p-Q} by $b'' \in R \setminus \{0\}$, we obtain equations with coefficients entirely in $R$:
\begin{align}
b' f &= q' m + r', && \text{with } b' \in R \setminus \{0\} \text{ and } q', r' \in R[x], \label{eq:cleared-f-div} \\
b'' p &= q'' m + r'', && \text{with } b'' \in R \setminus \{0\} \text{ and } q'', r'' \in R[x], \label{eq:cleared-p-div}
\end{align}
where $q' := b' \hat{q}$, $r' := b' \hat{r}$, $q'' := b'' \hat{\hat{q}}$, and $r'' := b'' \hat{\hat{r}}$ all lie in $R[x]$ by the choice of $b'$ and $b''$.
Multiplying by non-zero scalars in $R$ leaves polynomial degrees unchanged:
\[
r' = 0 \quad \text{or} \quad \deg(r') < \deg(m), \qquad \text{and} \qquad r'' = 0 \quad \text{or} \quad \deg(r'') < \deg(m).
\]
\end{content}

\begin{speech}[00:41:26 - 00:42:45]
When I clear the denominators, I don't change the degree here.

Now, this is in $I$, that is in $I$ because $m$ is in $I$, so $r'$ is in $I$. But $m$ is the least-degree element in $I$ --- you can't be lower than that.

So, these cases are not possible, and they must be zero. So, actually these terms go away. \inlinemetanote{erases $+ r'$ and $+ r''$}

So, we have $b' f = q' m$ \inlinemetanote{writes: $b' f = q' m$} and $b'' p = q'' m$. \inlinemetanote{writes: $b'' p = q'' m$}

And we can write this on the right-hand side as a content: this equals $q' c(m) m^*$ \inlinemetanote{writes: $= q' c(m) m^*$}, and this is $q'' c(m) m^*$. \inlinemetanote{writes: $= q'' c(m) m^*$}

So, that means $m^*$, which is primitive, divides this product.

Ah, I failed! \inlinemetanote{the bell rings, marking the end of the lecture hour}

I'll finish it next time. We're basically almost... I'll do it next time, there's no reason to go over. I need another three minutes.

This lecture was maybe a little bit faster than usual.

But... \inlinemetanote{the lecture ends; students applaud and pack up}
\end{speech}

\begin{content}[{Vanishing of Remainders and Divisibility by Primitive \texorpdfstring{$m^*$}{m*}}]
\begin{claim}[Vanishing of Remainders]\label[claim]{clm:remainders-zero}
The cleared remainders in \eqref{eq:cleared-f-div} and \eqref{eq:cleared-p-div} must vanish identically:
\[
r' = 0 \quad \text{and} \quad r'' = 0.
\]
\end{claim}

\begin{proof}[Proof of \cref{clm:remainders-zero}]
Rearranging \eqref{eq:cleared-f-div} and \eqref{eq:cleared-p-div} gives:
\begin{align*}
r' &= b' f - q' m, \\
r'' &= b'' p - q'' m.
\end{align*}
Recall that $I = (p, f)$. Since $f \in I$ and $m \in I$, any $R[x]$-linear combination of $f$ and $m$ belongs to $I$, so $r' \in I$.
Similarly, since $p \in I$ and $m \in I$, $r'' \in I$.

By definition \eqref{eq:m-minimal-def}, $m$ was chosen to be a non-zero element of minimal degree in $I$.
If $r' \neq 0$, then $r' \in I \setminus \{0\}$ with $\deg(r') < \deg(m)$, which contradicts the minimality of $\deg(m)$. Therefore, $r' = 0$.
The identical argument shows that $r'' = 0$.
\end{proof}

\begin{steps}[Divisibility of $b'f$ and $b''p$ by the Primitive Factor $m^*$]
Substituting $r' = 0$ and $r'' = 0$ into \eqref{eq:cleared-f-div} and \eqref{eq:cleared-p-div}, we obtain exact factorization identities in $R[x]$:
\begin{align}
b' f &= q' m, \label{eq:exact-bprime-f} \\
b'' p &= q'' m. \label{eq:exact-bprime-p}
\end{align}
Expressing $m$ via its canonical content decomposition $m = c(m) \cdot m^*$, where $c(m) \in R$ and $m^* \in R[x]$ is primitive, equations \eqref{eq:exact-bprime-f} and \eqref{eq:exact-bprime-p} become:
\begin{align}
b' f &= q' \cdot c(m) \cdot m^*, \label{eq:div-bprime-f-mstar} \\
b'' p &= q'' \cdot c(m) \cdot m^*. \label{eq:div-bprime-p-mstar}
\end{align}
Thus, the primitive polynomial $m^* \in R[x]$ satisfies:
\[
m^* \mid b' f \quad \text{and} \quad m^* \mid b'' p \quad \text{in } R[x],
\]
for non-zero scalars $b', b'' \in R \setminus \{0\}$.
By the divisibility property of primitive polynomials (\cref{clm:primitive-divisibility}), this immediately implies:
\[
m^* \mid f \quad \text{and} \quad m^* \mid p \quad \text{in } R[x].
\]
Here the lecture stops, at the bell. The proof that $\deg(m) = 0$, and from it $g \in (p)$, is continued in the next lecture; the remaining argument is given below.
\end{steps}
\end{content}

\begin{aiNote}[Beyond the Lecture: Completion of the Proof of Case 2]
% [PERSISTENT, ADDITION]
% The pipeline's "End of Lecture / Board State Checkpoint" note stood here. Resolved on 2026-10-24:
% the lecture ended normally at the bell, and the cut-off argument is completed below, as the
% lecturer announced for the next lecture. The checkpoint's board state is the last content box above.
The bell ended the lecture a few minutes before the end of the proof (\qt{I'll do it next time}). The remaining argument, which the next lecture is to give, runs as follows; it uses only the results of this lecture.
\begin{steps}[Completion of the Argument for $\deg(p) > 0$]
At the end of the lecture, $m^* \mid f$ and $m^* \mid p$ in $R[x]$, where $m^*$ is the primitive part of the element $m \in I = (p, f)$ of minimal degree.
\begin{itemize}
    \item \textbf{$m^*$ is a unit:} Write $p = m^* t$ with $t \in R[x]$. Since $p$ is irreducible, $m^*$ or $t$ is a unit. If $t$ were a unit, then $p = m^* t$ would divide $f$ (as $m^* \mid f$), that is, $f \in (p)$, contradicting the hypothesis $f \notin (p)$. Hence $m^*$ is a unit of $R[x]$, so $m^* \in R^\times$ has degree $0$ (\cref{thm:units-polynomial-ring-integral-domain}).
    \item \textbf{$m$ is a non-zero constant:} Since $m = c(m) \cdot m^*$ with both factors in $R$, $m =: c \in R \setminus \{0\}$; this is the claim $\deg(m) = 0$ written on the board.
    \item \textbf{$p$ is primitive:} $p = c(p) \cdot p^*$ with $\deg(p^*) = \deg(p) > 0$, so $p^*$ is not a unit, and the irreducibility of $p$ forces $c(p)$ to be a unit.
    \item \textbf{Deducing $g \in (p)$:} Because $c = m \in I = (p, f)$, there exist polynomials $A, B \in R[x]$ such that:
    \[
    c = A p + B f.
    \]
    Multiplying both sides by $g$ yields:
    \[
    c g = A p g + B (f g).
    \]
    Since $p \mid p g$ and by hypothesis $p \mid f g$, it follows that $p \mid cg$ in $R[x]$.
    Applying the divisibility property of primitive polynomials (\cref{clm:primitive-divisibility}) to the primitive polynomial $p$ and $p \mid c g$ with $c \in R \setminus \{0\}$, we conclude:
    \[
    p \mid g \quad \text{in } R[x] \implies g \in (p).
    \]
\end{itemize}
Together with Case~1, every irreducible element $p \in R[x]$ is prime (and $(p) \neq R[x]$, since $p$ is not a unit). By Step~1 of Theorem~A (\cref{clm:step1-uniqueness}), factorizations in $R[x]$ are therefore unique; with the existence proved in Step~1 above, this concludes the proof of Gauss's Theorem (\cref{thm:ufd-polynomial-ring-w6wed}).
\end{steps}
\end{aiNote}

\end{proofWrapper}

\begin{overview}[Summary of Key Concepts]
\begin{itemize}
    \item \textbf{Content of a Polynomial:} For a UFD $R$ and non-zero polynomial $f(x) = \sum_{j=0}^n a_j x^j \in R[x] \setminus \{0\}$, the content is $c(f) := \operatorname{GCD}(a_0, \dots, a_n) \in R / R^\times$, determined uniquely up to units in $R$.
    \item \textbf{Primitive Polynomials:} A polynomial $g \in R[x]$ is primitive if $c(g) \sim 1$. Every non-zero polynomial factors uniquely up to units as $f = c(f) \cdot f^*$, where $f^*$ is primitive.
    \item \textbf{Gauss's Lemma:} If $f, g \in R[x]$ are primitive, then their product $fg$ is primitive ($c(fg) \sim 1$). The proof utilizes reduction modulo a prime ideal $(p) \subset R$, exploiting the fact that $(R/(p))[x]$ is an integral domain.
    \item \textbf{Divisibility by Primitive Polynomials:} If $g \in R[x]$ is primitive and $g \mid bf$ for some non-zero scalar $b \in R$, then $g \mid f$ in $R[x]$.
    \item \textbf{Gauss's Theorem on Polynomial Rings:} If $R$ is a UFD, then $R[x]$ is a UFD:
    \begin{itemize}
        \item \emph{Existence of Factorization:} Proven by induction on polynomial degree. Factoring $f = c(f) f^*$ reduces the problem to primitive polynomials; reducibility $f^* = AB$ cannot yield degree-zero non-unit factors because $c(f^*) \sim 1$, ensuring strictly decreasing degrees $\deg(A), \deg(B) < \deg(f)$.
        \item \emph{Uniqueness of Factorization:} Proven by showing that every irreducible element is prime. Constant irreducibles ($\deg(p) = 0$) are handled by Gauss's Lemma and prime divisibility in $R$; positive-degree irreducibles ($\deg(p) > 0$) are analyzed using division with remainder in $Q[x] = (\operatorname{Frac}(R))[x]$ and clearing denominators. The bell ended the lecture just before the last step ($\deg(m) = 0$, hence $g \in (p)$), which the next lecture completes; the aiNote at the end of the proof gives it.
    \end{itemize}
\end{itemize}
\end{overview}